% Part: lambda-calculus
% Chapter: introduction
% Section: composition

\documentclass[../../../include/open-logic-section]{subfiles}

\begin{document}

\olfileid{lam}{int}{com}
\olsection{\usetoken{S}{lambda definable} વિધેયો સંયોજન હેઠળ બંધ છે}

\begin{lem}
!!{lambda definable} વિધેયો સંયોજન હેઠળ બંધ છે.
\end{lem}

\begin{proof}
ધારો કે $h$, $g_0,$ \dots,~$g_{k-1}$ના સંયોજનથી~$f$
વ્યાખ્યાયિત થાય છે. હવે ધારો કે $H$, $G_0$, \dots,~$G_{k-1}$ વડે
અનુક્રમે $h$, $g_0$, \dots,~$g_{k-1}$ !!{lambda defined} છે. આપણે
$f$ને !!{lambda define} કરતું પદ~$F$ શોધવાનું છે. પરંતુ~$F$ને
સીધું આ રીતે વ્યાખ્યાયિત કરી શકીએ:
\[
F(x_0, \dots, x_{l-1}) = H(G_0(x_0, \dots, x_{l-1}),
\dots, G_{k-1}(x_0, \dots, x_{l-1})).
\]
બીજા શબ્દોમાં, લૅમ્બડા કલનની ભાષા સંયોજનનું નિરૂપણ કરવા માટે
અનુકૂળ છે.
\end{proof}

\end{document}
